\section{GRPO：把 PPO 再简化一层}

\subsection{为什么需要 GRPO}

课堂给出的动机很直接：
\begin{itemize}
    \item PPO 太复杂，value model 很占资源。
    \item DPO 虽然简单，但它假设数据本身就是 pairwise preference。
    \item 现实里的 RLVR 场景往往不是纯 pairwise 数据，而是更一般的 rollout + verifiable reward。
\end{itemize}

于是就有了 \textbf{GRPO}。它的思路可以概括成一句话：\textbf{保留 PPO 的外壳，去掉 value function，把 advantage 变成 group 内标准化的分数。}

\begin{figure}[H]
\centering
\includegraphics[width=0.95\textwidth]{slides-images/slide-029.jpg}
\caption{为什么还需要新算法：PPO 太重，DPO 不一定适用}
\end{figure}

\begin{figure}[H]
\centering
\includegraphics[width=0.95\textwidth]{slides-images/slide-030.jpg}
\caption{GRPO：去掉 value function / advantage 估计，把优势定义成 group 内的 z-score}
\end{figure}

\subsection{GRPO 的简单实现}

GRPO 在代码里非常短。最核心的流程就是：
\begin{enumerate}
    \item 对同一个 prompt 采样一组 rollouts。
    \item 给每个 rollout 计算可验证 reward。
    \item 在 group 内做均值和方差归一化，得到 advantage。
    \item 加上 KL 约束，直接对 policy 做梯度更新。
\end{enumerate}

这也是为什么课堂里说它“看起来像是 policy gradient，但已经被整理得更像工程实现”。

\[
A_i = \frac{R_i - \mu_{\mathrm{group}}}{\sigma_{\mathrm{group}} + \epsilon}.
\]

如果把这个式子和 PPO 相比，你会发现它少了一个 value network，但代价是：这个“标准化”并不一定是严格无偏的 baseline。

\begin{figure}[H]
\centering
\includegraphics[width=0.95\textwidth]{slides-images/slide-031.jpg}
\caption{一个很短的 GRPO 代码实现：奖励、归一化、KL、梯度更新}
\end{figure}

\begin{figure}[H]
\centering
\includegraphics[width=0.95\textwidth]{slides-images/slide-032.jpg}
\caption{GRPO 的 advantage 计算：本质上就是 group 内的标准化}
\end{figure}

\begin{figure}[H]
\centering
\includegraphics[width=0.95\textwidth]{slides-images/slide-034.jpg}
\caption{GRPO 的一个关键问题：标准化项并不是一个严格无偏的 baseline}
\end{figure}

\subsection{length bias 和 unbiased 版本}

GRPO 的一个副作用是 \textbf{长度偏差}。因为它对 group 内的分数做归一化，模型可能会学到偏向更长或更短的回答。课堂里还提到了后续工作对这点的修正：如果希望更严格地保持梯度无偏，标准化项和长度归一化项都要重新审视。

\begin{figure}[H]
\centering
\includegraphics[width=0.95\textwidth]{slides-images/slide-035.jpg}
\caption{GRPO 的无偏版本讨论：标准化分数并不等于严格的 baseline}
\end{figure}

\begin{figure}[H]
\centering
\includegraphics[width=0.95\textwidth]{slides-images/slide-036.jpg}
\caption{GRPO 的长度偏差：分数归一化和长度归一化会改变模型偏好的回答长度}
\end{figure}

